\documentclass[11pt, a4paper]{article}
\usepackage[margin=1in]{geometry}
\usepackage{amsmath, amssymb, amsthm, mathtools}
\usepackage{hyperref}
\usepackage{graphicx}
\usepackage{booktabs}
\usepackage{enumitem}
\usepackage{listings}
\usepackage{xcolor}
\usepackage{cite}

\definecolor{codegreen}{rgb}{0,0.6,0}
\definecolor{codegray}{rgb}{0.5,0.5,0.5}
\definecolor{codepurple}{rgb}{0.58,0,0.82}
\definecolor{backcolour}{rgb}{0.95,0.95,0.92}

\lstdefinestyle{mystyle}{
    backgroundcolor=\color{backcolour},   
    commentstyle=\color{codegreen},
    keywordstyle=\color{magenta},
    numberstyle=\tiny\color{codegray},
    stringstyle=\color{codepurple},
    basicstyle=\ttfamily\footnotesize,
    breakatwhitespace=false,         
    breaklines=true,                 
    captionpos=b,                    
    keepspaces=true,                 
    numbers=left,                    
    numbersep=5pt,                  
    showspaces=false,                
    showstringspaces=false,
    showtabs=false,                  
    tabsize=2
}
\lstset{style=mystyle}

\newtheorem{theorem}{Theorem}
\newtheorem{lemma}[theorem]{Lemma}
\newtheorem{definition}{Definition}

\title{Prime-Indexed Dialectical Semantics\\
A Defensive Publication on PIRTM/DRMM,\\
Conceptual Contestation Fields, and Admissible Meaning Evolution}
\author{Defensive Research Dossier}
\date{\today}

\begin{document}
\maketitle

\begin{abstract}
This document consolidates and formalizes a framework for modeling conceptual evolution under strict epistemic and mathematical constraints.
It is intended as a defensive publication establishing prior art for Prime-Indexed Recursive Tensor Mathematics (PIRTM) and the Dialectical Robust Meaning Model (DRMM) as a \emph{certification calculus} over distributional semantics, not as a metaphysical theory of meaning.
We define the core objects (events, branches, fields), a prime-indexed multiplicity space, grounding coverage, a dialectical tension metric, a gated admissibility pipeline, and a minimal conformance protocol.
The central result is a protocol-level firewall: no raw vector or embedding trajectory may be treated as a concept or concept evolution path without passing explicit grounding, robustness, and dialectical non-collapse gates.
\end{abstract}
\newpage
\tableofcontents
\newpage
\section{Executive Summary}

\subsection{Motivation and Prior Art Gap}

Embedding-based distributional semantics has become a standard tool in natural language processing and AI.
However, most work implicitly collapses from:
\[
\text{Level 0 (embeddings)} \;\longrightarrow\; \text{Level 3 (concept evolution claims)}
\]
with no intervening semantics.
That is, a change in vector representation is taken as a proxy for conceptual change.

Philosophers of language, formal semanticists, and skeptical AI practitioners have repeatedly noted that this is a category error:
\begin{itemize}[nosep]
  \item Distributional models capture usage and association, not full truth-conditional content.
  \item Text-only models are ungrounded and cannot, by themselves, secure reference.
  \item Embedding spaces blur relation types (synonymy, antonymy, rhetorical opposition, institutional coupling, etc.).
  \item Diachronic embedding studies are sensitive to corpus and model choices, making naive concept-drift claims unstable.
\end{itemize}

This defensive publication addresses that gap by:
\begin{enumerate}[nosep]
  \item Defining PIRTM: a prime-indexed multiplicity framework where events, branches, and fields live in a mathematically structured space.
  \item Defining DRMM: a dialectical update and certification operator on top of distributional semantics, with explicit gate operators and rejection semantics.
  \item Providing a formal admissibility specification (Appendix-style) and a minimal evaluation protocol that together turn DRMM into a testable, reproducible protocol.
\end{enumerate}

\subsection{Core Idea: Conceptual Contestation Fields}

Instead of modeling ``concepts themselves'', PIRTM/DRMM models \emph{conceptual contestation fields}:
\begin{quote}
Temporally evolving distributions of usage, stance, and institutional uptake, constrained by distributional evidence, dialectical collisions, and explicit grounding and provenance.
\end{quote}

A ``concept'' in this framework is not a single vector or latent variable, but a recursively stabilized pattern across:
\begin{itemize}[nosep]
  \item text (linguistic usage),
  \item practice (experiments, code, data),
  \item institution (regulations, definitions, adoption decisions).
\end{itemize}

PIRTM provides the prime-indexed multiplicity space; DRMM is the evolution and certification operator that enforces epistemic discipline.

\subsection{Key Invariants (Semantic Safety Firewall)}

The framework is characterized by the following non-negotiable invariants:
\begin{enumerate}[nosep]
  \item \textbf{No direct Level 0 $\to$ Level 3 mapping.} There is no admissible morphism from raw embeddings to concept evolution claims.
  \item \textbf{Gate-based certification.} Any ``candidate conceptual evolution'' must pass three gates:
    \begin{itemize}[nosep]
      \item Grounding gate $\Pi_G$ (sufficient extra-linguistic evidence),
      \item Robustness gate $\Pi_R$ (stability under perturbations),
      \item Dialectical gate $\Pi_D$ (non-collapse of high-tension branches).
    \end{itemize}
  \item \textbf{Typed rejection states.} Failures are reported as typed rejections ($\bot_G,\bot_R,\bot_D$), not silently dropped.
  \item \textbf{Pluralism preserved.} High-tension, well-grounded branches cannot be averaged into a single ``concept'' trajectory.
  \item \textbf{Provenance mandatory.} Every output must declare corpus, model, clustering, perturbations, metrics, thresholds, and grounding model.
\end{enumerate}

Together, these make DRMM a \emph{semantic safety specification} on top of distributional semantics.

\section{Comprehensive Mathematical Overview}

\subsection{PIRTM Multiplicity Space}

\subsubsection*{Prime-Indexed State Space}

Prime-Indexed Recursive Tensor Mathematics (PIRTM) assumes a universe of primes $\mathbb{P}$ indexing irreducible interaction types.
For semantic modeling, we consider a Hilbert-like feature space
\[
\mathcal{H}
  \;=\;
  \ell^2(\mathbb{P})
  \,\widehat{\otimes}\,
  L^2(\mathbb{R})
  \,\widehat{\otimes}\,
  \mathbb{R}^d
  \,\widehat{\otimes}\,
  \mathcal{H}_{\text{stance}}
  \,\widehat{\otimes}\,
  \mathcal{H}_{\text{inst}},
\]
where:
\begin{itemize}[nosep]
  \item $\ell^2(\mathbb{P})$ encodes prime-labeled event types $p$,
  \item $L^2(\mathbb{R})$ encodes time $t$,
  \item $\mathbb{R}^d$ encodes distributional embeddings,
  \item $\mathcal{H}_{\text{stance}}$ encodes stance/affect coordinates,
  \item $\mathcal{H}_{\text{inst}}$ encodes institutional features.
\end{itemize}

A PIRTM \emph{event} is a vector in $\mathcal{H}$ enriched with discrete metadata (actors, contexts, grounding hooks).

\subsubsection*{Events}

Each elementary event $e$ is a tuple:
\[
e
=
\langle
  p,\;
  t,\;
  c,\;
  a,\;
  G(e),\;
  r
\rangle
\]
with:
\begin{itemize}[nosep]
  \item $p \in \mathbb{P}$: prime / interaction type,
  \item $t \in \mathbb{T}$: time index,
  \item $c \in \mathcal{C}$: context (corpus slice, community, domain),
  \item $a \in \mathcal{A}$: actor/institution,
  \item $G(e) \subseteq \mathcal{G}$: grounding hooks,
  \item $r \in \mathbb{R}^d$: embedding of the usage.
\end{itemize}

Grounding hooks are typed:
\[
\mathcal{T}_G = \{\tau_1,\dots,\tau_M\}, \quad
\mathcal{G} = \bigsqcup_{\tau \in \mathcal{T}_G} \mathcal{G}_\tau, \quad
\mathrm{type}(g) \in \mathcal{T}_G.
\]

\subsection{Branches, Fields, and Grounding Coverage}

\subsubsection*{Branch States}

Events for a term $\omega$ are partitioned (via clustering $\mathcal{C}_\omega$) into branch candidates.
For branch $j$ at time $t$ with event set $V_j(t)$, define:
\[
\theta_j(t) = \bigl(\mu_j(t),\, \sigma_j(t),\, \iota_j(t),\, \gamma_j(t)\bigr),
\]
where:
\begin{itemize}[nosep]
  \item $\mu_j(t) \in \mathbb{R}^d$: aggregated embedding (e.g.\ mean in $V_j$),
  \item $\sigma_j(t)$: stance/affect encoding (e.g.\ from sentiment/stance models),
  \item $\iota_j(t)$: institutional encoding (vector or discrete profile),
  \item $\gamma_j(t) \in [0,1]$: grounding coverage.
\end{itemize}

\subsubsection*{Grounding Coverage $\gamma_j$}

For branch $j$, aggregate hooks:
\[
\mathcal{G}_j = \biguplus_{e \in V_j} G(e).
\]
For each type $\tau \in \mathcal{T}_G$:
\[
n_j(\tau) = \bigl|\{g \in \mathcal{G}_j : \mathrm{type}(g) = \tau\}\bigr|, \quad
1_j(\tau) = \mathbf{1}[n_j(\tau) > 0].
\]
\textbf{Type coverage}:
\[
C_j^{\mathrm{type}} = \frac{\sum_{\tau \in \mathcal{T}_G} 1_j(\tau)}{|\mathcal{T}_G|}.
\]
Let $w_\tau \ge 0$, $\sum_\tau w_\tau = 1$, and saturation constants $K_\tau \ge 1$.
Define:
\[
S_j(\tau) = \min\!\left(1,\frac{n_j(\tau)}{K_\tau}\right),
\qquad
C_j^{\mathrm{strength}} = \sum_{\tau \in \mathcal{T}_G} w_\tau\, S_j(\tau).
\]
Finally:
\[
\gamma_j = \lambda\, C_j^{\mathrm{type}} + (1-\lambda)\, C_j^{\mathrm{strength}},
\quad \lambda \in [0,1],
\]
so that $\gamma_j \in [0,1]$ and encodes both diversity and density of grounding.

\subsubsection*{Fields and Tension Metric}

The \emph{field} for term $\omega$ at time $t$ is:
\[
\mathcal{F}_\omega(t) = \bigl(\{\theta_j(t)\}_j,\; \delta_{jk}(t)\bigr),
\]
where $\delta_{jk}(t) = \delta(\theta_j(t),\theta_k(t))$ is the dialectical tension metric.

Let:
\begin{itemize}[nosep]
  \item $\mu_i,\mu_j \in \mathbb{R}^d$,
  \item $\sigma_i,\sigma_j$ stance encodings,
  \item $\iota_i,\iota_j$ institutional encodings,
  \item $\gamma_i,\gamma_j \in [0,1]$ grounding coverage.
\end{itemize}
With component distances $d_{\mathrm{stance}}, d_{\mathrm{inst}} \ge 0$ and weights $\alpha,\beta,\gamma,\eta \ge 0$:
\[
\delta(\theta_i,\theta_j)
=
\alpha \,\lVert \mu_i-\mu_j \rVert_2
+
\beta \, d_{\mathrm{stance}}(\sigma_i,\sigma_j)
+
\gamma \, d_{\mathrm{inst}}(\iota_i,\iota_j)
+
\eta \, |\gamma_i-\gamma_j|.
\]

This emphasizes asymmetry of grounding: branches with very different $\gamma$ incur additional tension.

\subsection{DRMM as a Certification Operator}

\subsubsection*{Pipeline}

For each term $\omega$:
\[
\Pi_{\mathrm{DRMM}}(\omega)
=
\Pi_D \circ \Pi_R \circ \Pi_G \circ \mathcal{C}_\omega.
\]

\begin{itemize}[nosep]
  \item $\mathcal{C}_\omega$: aggregates events into branch candidates and computes $\theta_j$.
  \item $\Pi_G$: filters branches by grounding.
  \item $\Pi_R$: filters branches by robustness under perturbations.
  \item $\Pi_D$: enforces non-collapse and assigns field-level status.
\end{itemize}

No direct mapping from Level 0 (embeddings) to Level 3 (conceptual evolution) is permitted; there is no admissible $\Pi_{03}$.

\subsubsection*{Gates}

\paragraph{Grounding gate $\Pi_G$.}
For threshold $\gamma_{\min}$:
\[
\Pi_G(\theta_j) =
\begin{cases}
\theta_j & \gamma_j \ge \gamma_{\min},\\[2pt]
\bot_G   & \text{otherwise.}
\end{cases}
\]

\paragraph{Robustness gate $\Pi_R$.}
Given perturbation runs $k$ and $d_{\mathrm{branch}}$:
\[
\Pi_R(\theta_j) =
\begin{cases}
\theta_j & \max_k d_{\mathrm{branch}}(\theta_j, \theta_j^{(k)}) < \epsilon_R,\\[2pt]
\bot_R   & \text{otherwise.}
\end{cases}
\]

\paragraph{Dialectical gate $\Pi_D$.}
Given threshold $\delta_{\mathrm{thr}}$, define high-tension graph $H$ with edges $(j,k)$ where $\delta_{jk} \ge \delta_{\mathrm{thr}}$.
$\Pi_D$ must:
\begin{itemize}[nosep]
  \item prohibit merging branches across any edge in $H$ (non-collapse),
  \item label fields with \texttt{persistent\_pluralism} when multiple grounded, robust, high-tension branches persist,
  \item allow \texttt{candidate\_evolution} only when a branch becomes singular and stable within a formerly contested component.
\end{itemize}

\subsubsection*{Branch Persistence}

Identity tracking uses $\delta$ and a threshold $\tau_{\mathrm{id}}$:
\[
\mathrm{Track}(\theta_j(t)) =
\begin{cases}
k^* & k^* = \arg\min_k \delta(\theta_j(t), \theta_k(t+\Delta t)) \ \text{and}\ \delta(\theta_j(t),\theta_{k^*}(t+\Delta t)) < \tau_{\mathrm{id}},\\[2pt]
\text{none} & \text{otherwise.}
\end{cases}
\]

This forbids informal storytelling about branch continuity; identity must be justified via tension.

\subsection{Output Schema and Rejection Semantics}

A certified output for term $\omega$ has shape:
\[
\mathcal{E}_\omega =
\bigl(
  \{\theta_j\}_j,\;
  \delta_{jk},\;
  \mathrm{status},\;
  \mathrm{provenance},\;
  \mathrm{thresholds}
\bigr),
\]
with:
\begin{itemize}[nosep]
  \item branch statuses in
  \[
  \{\texttt{candidate\_evolution},\ \texttt{persistent\_pluralism},\ \bot_G,\ \bot_R,\ \bot_D\},
  \]
  \item no raw Level 0 vectors or ``concept\_vector'' fields,
  \item complete provenance (corpus, time windows, model, clustering, perturbations, metrics, grounding model, thresholds).
\end{itemize}

Rejection states are part of the public output, not hidden:
\[
\bot_G: \text{insufficient grounding}, \quad
\bot_R: \text{insufficient robustness}, \quad
\bot_D: \text{dialectical constraint violation}.
\]
\newpage
\section{Minimal Evaluation Protocol (MEP-1)}

To make DRMM falsifiable and reproducible, we define a Minimal Evaluation Protocol (MEP-1) that exercises all normative clauses.

\subsection{Corpus Requirements}

MEP-1 requires at least two corpora:
\begin{itemize}[nosep]
  \item a \textbf{closed conceptual change} corpus (e.g.\ Pluto/``planet'' reclassification),
  \item a \textbf{contested live discourse} corpus (e.g.\ AGI/alignment/safety).
\end{itemize}

Each must:
\begin{enumerate}[nosep]
  \item Provide temporal metadata for windows (e.g.\ pre/post institutional act).
  \item Include actor/institution signal (authors, venues, organizations).
  \item Support extraction of grounding hooks (experiments, datasets, code, regulations, KB entities).
  \item Have sufficient event density per window for branch extraction.
\end{enumerate}

\subsection{Required Perturbations}

Robustness must be tested via at least three perturbation families:
\begin{itemize}[nosep]
  \item seed perturbation (multiple random seeds),
  \item corpus perturbation (bootstrap or subsampling),
  \item temporal perturbation (slightly shifted windows).
\end{itemize}
Embedding perturbation (alternate encoder) is recommended.
For each branch, $\theta_j^{(k)}$ must be computed under these perturbations.

\subsection{Mandatory Gate Tests}

MEP-1 specifies scenarios with expected outcomes:
\begin{itemize}[nosep]
  \item \textbf{Grounding rejection}: an ungrounded rhetorical branch triggers $\bot_G$.
  \item \textbf{Robustness rejection}: an unstable branch across perturbations triggers $\bot_R$.
  \item \textbf{Persistent pluralism}: multiple grounded, robust branches with high $\delta$ yield \texttt{persistent\_pluralism}.
  \item \textbf{Candidate evolution}: in the Pluto/planet case, a scientific branch that becomes dominant and stable yields \texttt{candidate\_evolution}.
\end{itemize}

\subsection{Branch Persistence and Non-Collapse Tests}

Identity tracking and non-collapse are tested by:
\begin{itemize}[nosep]
  \item continuity cases (small $\delta$ across windows produces consistent tracking),
  \item branch death cases (no suitable successor; identity terminates),
  \item high-tension pairs (implementations must not merge them).
\end{itemize}

A single merged trajectory for branches with $\delta\ge\delta_{\mathrm{thr}}$ constitutes a protocol-level failure.

\subsection{Provenance Conformance and Output Validation}

MEP-1 requires:
\begin{itemize}[nosep]
  \item machine-readable provenance with declared thresholds and grounding model,
  \item no hidden branch suppression or unreported failures,
  \item no Level 0 exposure in the public schema,
  \item statuses aligned with gate behavior.
\end{itemize}

\section{Illustrative Code Snippet}

The following pseudo-code illustrates how an implementation might instantiate the grounding coverage and gate composition.
This is not normative code, but a reference sketch.

\subsection*{Grounding Coverage $\gamma_j$}

\begin{verbatim}
def compute_gamma(G_j, types, weights, saturation, lam):
    # G_j: list of grounding hooks for branch j
    # types: list of grounding types tau
    # weights: dict tau -> w_tau
    # saturation: dict tau -> K_tau
    # lam: lambda in [0, 1]
    # 1. Count hooks per type
    n = {tau: 0 for tau in types}
    for g in G_j:
        tau = g.type
        if tau in n:
            n[tau] += 1
    # 2. Type coverage
    type_hits = sum(1 for tau in types if n[tau] > 0)
    C_type = type_hits / max(1, len(types))
    # 3. Strength coverage
    S = {}
    for tau in types:
        K_tau = saturation.get(tau, 1)
        S[tau] = min(1.0, n[tau] / max(1, K_tau))
    C_strength = sum(weights.get(tau, 0.0) * S[tau] for tau in types)
    # 4. Combined coverage
    gamma = lam * C_type + (1.0 - lam) * C_strength
    return max(0.0, min(1.0, gamma))
\end{verbatim}

\subsection*{Gate Composition and Status Assignment}

\begin{verbatim}
def certify_branch(theta_j, perturbations, thresholds):
    gamma_min = thresholds["gamma_min"]
    eps_R = thresholds["epsilon_R"]
    # Grounding gate
    if theta_j.gamma < gamma_min:
        return ("⊥G", theta_j)
    # Robustness gate
    d_max = 0.0
    for theta_j_k in perturbations:
        d = branch_distance(theta_j, theta_j_k)
        d_max = max(d_max, d)
    if d_max >= eps_R:
        return ("⊥R", theta_j)
    # Passed G and R; D is applied at field level
    return ("ok", theta_j)

def certify_field(field, thresholds):
    # field: (branches, delta_matrix)
    branches = field.branches
    delta = field.delta
    delta_thr = thresholds["delta_thr"]
    # Build high-tension graph
    H = build_graph(branches, delta, delta_thr)
    # Non-collapse enforced by not merging connected components in H
    # Determine field-level status
    grounded_robust = [b for b in branches if b.status == "ok"]
    high_tension_components = connected_components(H)
    if any(len(comp) >= 2 for comp in high_tension_components):
        return "persistent_pluralism"
    elif len(grounded_robust) == 1:
        return "candidate_evolution"
    else:
        return "undefined_or_trivial"
\end{verbatim}

This code sketch demonstrates:
\begin{itemize}[nosep]
  \item explicit $\gamma_j$ computation from typed hooks,
  \item gate-wise statuses ($\bot_G$, $\bot_R$, and ``ok''),
  \item field-level dialectical status without any averaging across high-tension branches.
\end{itemize}

\section{Conclusion and Defensive Claims}

This report establishes:
\begin{enumerate}[nosep]
  \item A prime-indexed multiplicity framework (PIRTM) for events and semantic fields.
  \item A dialectical certification operator (DRMM) that:
    \begin{itemize}[nosep]
      \item forbids direct vector $\to$ concept mappings,
      \item imposes three epistemic gates (grounding, robustness, dialectical non-collapse),
      \item uses typed rejection states and persistent pluralism as first-class outcomes,
      \item requires machine-readable provenance for reproducibility.
    \end{itemize}
  \item A minimal evaluation protocol (MEP-1) that operationalizes the spec and makes it testable.
\end{enumerate}

The key defensive contribution is the transformation of embedding-based ``semantic change'' from an informal metaphor into a constrained reporting protocol for conceptual contestation, with clear admissibility, rejection, and non-collapse semantics.
Any future work that:
\begin{itemize}[nosep]
  \item uses prime-indexed multiplicity spaces to model usage, stance, and institutional uptake,
  \item passes vector-based meaning hypotheses through explicit grounding, robustness, and dialectical gates,
  \item and encodes pluralism and rejection as explicit statuses,
\end{itemize}
will necessarily intersect with the prior art established here.
\newpage
\appendix
\section*{Appendix A \quad DRMM Admissibility Specification}

This appendix specifies the admissibility conditions for Dialectical Robust Meaning Model (DRMM) outputs. It is written as a protocol: it defines the typed objects, gate operators, rejection states, identity rules, tension metric, output schema, provenance requirements, and forbidden operations that any implementation must satisfy to claim DRMM compliance.

Throughout, we distinguish a thin human–readable gloss (in prose) from a canonical operational form (in equations).

\subsection*{A.1 \quad Admissibility Semantics}

\paragraph{Human–readable.}
A DRMM output for a term $\omega$ is admissible only if it is produced by the certified pipeline; no direct promotion from event–level embeddings (Level 0) to ``concept evolution'' (Level 3) is allowed.

\paragraph{Canonical form.}
For a target term $\omega$, define the certified operator
\[
\Pi_{\mathrm{DRMM}}(\omega) \;=\; \Pi_D \circ \Pi_R \circ \Pi_G \circ \mathcal{C}_\omega,
\]
where $\mathcal{C}_\omega$ is a branch constructor (Section~\ref{sec:A2}), and $\Pi_G, \Pi_R, \Pi_D$ are the grounding, robustness, and dialectical gates (Section~\ref{sec:A3}).

An output $\mathcal{E}_\omega$ is \emph{DRMM–admissible} iff there exists a run of this composition that produces $\mathcal{E}_\omega$ under a declared provenance configuration (Section~\ref{sec:A8}).

We explicitly forbid any direct morphism
\[
\Pi_{03} : \text{Level 0} \to \text{Level 3}
\]
in the admissible calculus. Informally: \emph{no raw Level–0 vector is a reportable semantic object, and no \texttt{concept\_vector} field is admissible as evidence of conceptual evolution.}

\subsection*{A.2 \quad Typed Objects and Branch Space}
\label{sec:A2}

\paragraph{Human–readable.}
DRMM operates over typed events, branches, fields, and statuses, not bare vectors.

\paragraph{Canonical form.}

\subsubsection*{A.2.1 \quad Events and Grounding Hooks}

Fix a finite set of grounding types
\[
\mathcal{T}_G = \{\tau_1,\dots,\tau_M\},
\]
e.g.\ $\{\texttt{experiment}, \texttt{dataset}, \texttt{code}, \texttt{regulation}, \texttt{kb\_entity}\}$.

Let
\[
\mathcal{G} \;=\; \bigsqcup_{\tau \in \mathcal{T}_G} \mathcal{G}_\tau
\]
be the disjoint union of identifiers of each type $\tau$.

Each event $e$ is a tuple
\[
e \;=\; \langle p,\; t,\; c,\; a,\; G(e),\; r \rangle
\]
with:
\begin{itemize}
  \item $p \in \mathbb{P}$: prime / interaction type,
  \item $t \in \mathbb{T}$: time index,
  \item $c \in \mathcal{C}$: context ID (e.g.\ corpus slice or community),
  \item $a \in \mathcal{A}$: actor or institution ID,
  \item $G(e) \subseteq \mathcal{G}$: set of grounding hooks for $e$,
  \item $r \in \mathbb{R}^d$: event–level embedding (Level 0).
\end{itemize}
Each hook $g \in \mathcal{G}$ has a declared type $\mathrm{type}(g) \in \mathcal{T}_G$.

\subsubsection*{A.2.2 \quad Branch States}

A branch $j$ at time $t$ is represented as
\[
\theta_j(t) = \bigl(\mu_j(t),\, \sigma_j(t),\, \iota_j(t),\, \gamma_j(t)\bigr),
\]
where:
\begin{itemize}
  \item $\mu_j(t) \in \mathbb{R}^d$: branch embedding (aggregated geometry),
  \item $\sigma_j(t)$: stance / affect profile,
  \item $\iota_j(t)$: institutional footprint,
  \item $\gamma_j(t) \in [0,1]$: grounding coverage (defined in Section~\ref{sec:A2gamma}).
\end{itemize}
Let $V_j(t)$ be the set of events assigned to branch $j$ at time $t$.

\subsubsection*{A.2.3 \quad Grounding Coverage}
\label{sec:A2gamma}

\paragraph{Human–readable.}
$\gamma_j$ is a normalized score in $[0,1]$ that increases as a branch accumulates more and more diverse extra–linguistic evidence (experiments, datasets, code, regulations, entities), and saturates once each evidence type is sufficiently represented.

\paragraph{Canonical form.}
For branch $j$, define the multiset of hooks across its events:
\[
\mathcal{G}_j = \biguplus_{e \in V_j} G(e).
\]

For each grounding type $\tau \in \mathcal{T}_G$, let:
\[
n_j(\tau) = \bigl|\{g \in \mathcal{G}_j \mid \mathrm{type}(g)=\tau\}\bigr|, \quad
1_j(\tau) = \mathbf{1}\bigl[n_j(\tau) > 0\bigr].
\]

\emph{Type coverage} is
\[
C_j^{\mathrm{type}} = \frac{\sum_{\tau \in \mathcal{T}_G} 1_j(\tau)}{|\mathcal{T}_G|}.
\]

Let $w_\tau \ge 0$ be type weights with $\sum_{\tau} w_\tau = 1$, and $K_\tau \ge 1$ be saturation constants. Define
\[
S_j(\tau) = \min\!\left(1, \frac{n_j(\tau)}{K_\tau}\right),
\]
and the \emph{strength coverage}
\[
C_j^{\mathrm{strength}} = \sum_{\tau \in \mathcal{T}_G} w_\tau\, S_j(\tau).
\]

For $\lambda \in [0,1]$, define combined grounding coverage
\[
\gamma_j = \lambda\, C_j^{\mathrm{type}} + (1-\lambda)\, C_j^{\mathrm{strength}} \in [0,1].
\]

The tuple $(\mathcal{T}_G, w_\tau, K_\tau, \lambda)$ is implementation–configurable but must be fully declared in provenance (Section~\ref{sec:A8}).

\subsubsection*{A.2.4 \quad Fields and Statuses}

The \emph{field} for $\omega$ at time $t$ is
\[
\mathcal{F}_\omega(t) = \bigl(\{\theta_j(t)\}_j,\; \delta_{jk}(t)\bigr),
\]
where $\delta_{jk}(t) = \delta(\theta_j(t), \theta_k(t))$ is the pairwise tension metric (Section~\ref{sec:A6}).

Every branch receives a status from the finite set
\[
\mathrm{status}_j \in 
\{
\texttt{candidate\_evolution},\ 
\texttt{persistent\_pluralism},\ 
\bot_G,\ 
\bot_R,\ 
\bot_D
\}.
\]

\subsection*{A.3 \quad Gate Operators}
\label{sec:A3}

Each gate has a human description and a canonical admissibility rule.

\subsubsection*{A.3.1 \quad Grounding Gate $\Pi_G$}

\paragraph{Human–readable.}
Reject branches that lack sufficient extra–linguistic grounding.

\paragraph{Canonical form.}
Given $\gamma_{\min} \in (0,1]$, define
\[
\Pi_G(\theta_j) =
\begin{cases}
\theta_j & \text{if } \gamma_j \ge \gamma_{\min},\\[4pt]
\bot_G   & \text{otherwise.}
\end{cases}
\]

\subsubsection*{A.3.2 \quad Robustness Gate $\Pi_R$}

\paragraph{Human–readable.}
Reject branches whose trajectories are unstable under declared perturbations (corpus, seed, alignment, etc.).

\paragraph{Canonical form.}
Let $\theta_j^{(k)}$ be branch states obtained under perturbation runs $k=1,\dots,K$. Let $d_{\mathrm{branch}}$ be the approved branch distance metric. For robustness threshold $\epsilon_R > 0$,
\[
\Pi_R(\theta_j) =
\begin{cases}
\theta_j & \text{if } \max_k d_{\mathrm{branch}}(\theta_j, \theta_j^{(k)}) < \epsilon_R,\\[4pt]
\bot_R   & \text{otherwise.}
\end{cases}
\]
The construction of $d_{\mathrm{branch}}$ must be declared in provenance.

\subsubsection*{A.3.3 \quad Dialectical Gate $\Pi_D$}

\paragraph{Human–readable.}
Prevent silent averaging across high–tension branches; represent unresolved conflict as pluralism.

\paragraph{Canonical form.}
Given a field $\mathcal{F}_\omega = (\{\theta_j\}, \delta_{jk})$ and tension threshold $\delta_{\mathrm{thr}} > 0$, define the high–tension graph $H$ with vertex set $\{j\}$ and edges
\[
(j,k) \in E(H) \iff \delta_{jk} \ge \delta_{\mathrm{thr}}.
\]

$\Pi_D$ must satisfy:

\begin{enumerate}[label=(\roman*)]
  \item \textbf{Non–collapse.} For any edge $(j,k) \in H$, $\Pi_D$ shall not produce a single merged branch replacing both $\theta_j$ and $\theta_k$.
  \item \textbf{Status assignment.}
  \begin{itemize}
    \item If at least one branch survives all gates and all connected components of $H$ have size $\ge 2$, then the field–level status is \texttt{persistent\_pluralism}.
    \item If, over time, a connected component of $H$ diminishes (e.g.\ branches lose grounding or robustness) leaving a single surviving branch for that component, that branch may be eligible for \texttt{candidate\_evolution} (subject to Section~\ref{sec:A7}).
  \end{itemize}
\end{enumerate}

Branches that fail $\Pi_G$ or $\Pi_R$ already carry $\bot_G$ or $\bot_R$; $\Pi_D$ must not mask these failures.

\subsection*{A.4 \quad Rejection States}

\paragraph{Human–readable.}
Failures are typed outputs, not internal exceptions.

\paragraph{Canonical form.}
We define three rejection states:
\[
\bot_G \quad \text{(insufficient grounding)}, \qquad
\bot_R \quad \text{(insufficient robustness)}, \qquad
\bot_D \quad \text{(dialectical constraint violation)}.
\]
Every reported branch carries one of the statuses in A.2.4; silent dropping of failed branches is disallowed.

\subsection*{A.5 \quad Branch Persistence}
\label{sec:A5}

\paragraph{Human–readable.}
We require an explicit rule for when branch $j$ at $t_1$ is treated as the same branch at $t_2$.

\paragraph{Canonical form.}
Given time steps $t$ and $t+\Delta t$, branches $\{\theta_j(t)\}$ and $\{\theta_k(t+\Delta t)\}$, and the tension metric $\delta$:

\[
\mathrm{Track}(\theta_j(t)) =
\begin{cases}
k^* & \text{if } k^* = \arg\min_k \delta(\theta_j(t), \theta_k(t+\Delta t))\\
    & \quad \text{and } \delta(\theta_j(t), \theta_{k^*}(t+\Delta t)) < \tau_{\mathrm{id}},\\[4pt]
\text{none} & \text{otherwise,}
\end{cases}
\]
with identity threshold $\tau_{\mathrm{id}} > 0$.

If $\mathrm{Track} = \text{none}$, branch $\theta_j$ is treated as terminated or split; identity must not be silently preserved past this threshold.

\subsection*{A.6 \quad Tension Metric $\delta$}
\label{sec:A6}

\paragraph{Human–readable.}
Tension measures how far apart two branches are in geometry, stance, institution, and grounding asymmetry.

\paragraph{Canonical form.}
Let
\begin{itemize}
  \item $\mu_i,\mu_j \in \mathbb{R}^d$: branch embeddings,
  \item $\sigma_i,\sigma_j$: stance / affect encodings,
  \item $\iota_i,\iota_j$: institutional encodings,
  \item $\gamma_i,\gamma_j \in [0,1]$: grounding coverage.
\end{itemize}

Let $d_{\mathrm{stance}}(\sigma_i,\sigma_j) \ge 0$ and $d_{\mathrm{inst}}(\iota_i,\iota_j) \ge 0$ be component distances, and let $\alpha,\beta,\gamma,\eta \ge 0$ be weights. Define
\[
\delta(\theta_i,\theta_j)
=
\alpha \,\lVert \mu_i-\mu_j \rVert_2
+
\beta \, d_{\mathrm{stance}}(\sigma_i,\sigma_j)
+
\gamma \, d_{\mathrm{inst}}(\iota_i,\iota_j)
+
\eta \, |\gamma_i-\gamma_j|.
\]
The choice of component metrics and weights must be declared in provenance.

\subsection*{A.7 \quad Output Schema}

\paragraph{Human–readable.}
Every certified output must be structurally comparable across runs and implementations.

\paragraph{Canonical form.}
For term $\omega$, an admissible DRMM output has the shape
\[
\mathcal{E}_\omega = 
\bigl(
\{\theta_j\}_j,\;
\delta_{jk},\;
\mathrm{status},\;
\mathrm{provenance},\;
\mathrm{thresholds}
\bigr),
\]
where:
\begin{itemize}
  \item $\{\theta_j\}_j$ are branch states that have passed $\Pi_G$ and $\Pi_R$ (per–branch statuses still carried),
  \item $\delta_{jk} = \delta(\theta_j,\theta_k)$ for all reported pairs,
  \item $\mathrm{status}$ is a field–level status (e.g.\ \texttt{candidate\_evolution} or \texttt{persistent\_pluralism}),
  \item $\mathrm{provenance}$ is as defined in Section~\ref{sec:A8},
  \item $\mathrm{thresholds}$ collects the numeric values of $\gamma_{\min}, \epsilon_R, \delta_{\mathrm{thr}}, \tau_{\mathrm{id}}$.
\end{itemize}

Schema–level constraint: no field may expose raw Level–0 vectors or derived ``concept vectors'' as standalone semantic objects. Embeddings may appear only as components of $\theta_j$ (e.g.\ $\mu_j$) with clear typing.

\subsection*{A.8 \quad Provenance Requirements}
\label{sec:A8}

\paragraph{Human–readable.}
Two DRMM outputs are only comparable if we know how they were produced.

\paragraph{Canonical form.}
At minimum, $\mathrm{provenance}$ must include:
\[
\mathrm{provenance} =
(
\texttt{corpus\_id},
\texttt{time\_window},
\texttt{embedding\_model},
\texttt{clustering\_method},
\texttt{perturbation\_family},
\texttt{branch\_metric},
\texttt{grounding\_model},
\texttt{threshold\_config}
),
\]
where each component is a structured record:
\begin{itemize}
  \item \texttt{corpus\_id}: identifier(s) of source corpora / domains,
  \item \texttt{time\_window}: definition of temporal slices,
  \item \texttt{embedding\_model}: model name, version, key hyperparameters,
  \item \texttt{clustering\_method}: branch constructor method and parameters,
  \item \texttt{perturbation\_family}: types and counts of perturbations used in $\Pi_R$,
  \item \texttt{branch\_metric}: instantiation of $d_{\mathrm{branch}}$ and component distances,
  \item \texttt{grounding\_model}: the tuple $(\mathcal{T}_G, w_\tau, K_\tau, \lambda)$ and hook–extraction strategy,
  \item \texttt{threshold\_config}: values of $\gamma_{\min}, \epsilon_R, \delta_{\mathrm{thr}}, \tau_{\mathrm{id}}$.
\end{itemize}
No DRMM output is admissible without a complete provenance record.

\subsection*{A.9 \quad Normative vs.\ Configurable Parameters}

\paragraph{Human–readable.}
We separate framework invariants from tunable implementation choices.

\paragraph{Canonical form.}

\textbf{Normative (required for DRMM compliance):}
\begin{itemize}
  \item No Level–0 $\to$ Level–3 bypass (no $\Pi_{03}$).
  \item Existence and use of $\Pi_G, \Pi_R, \Pi_D$ as defined.
  \item Typed rejection states $\bot_G, \bot_R, \bot_D$.
  \item Branch identity tracking via $\delta$ and $\tau_{\mathrm{id}}$ (Section~\ref{sec:A5}).
  \item High–$\delta$ non–collapse constraint in $\Pi_D$.
  \item Provenance reporting as in Section~\ref{sec:A8}.
  \item Output schema shape in Section~\ref{sec:A7}.
\end{itemize}

\textbf{Configurable (tunable but must be declared):}
\begin{itemize}
  \item Threshold values: $\gamma_{\min}, \epsilon_R, \delta_{\mathrm{thr}}, \tau_{\mathrm{id}}$.
  \item Choice of clustering algorithm and parameters in $\mathcal{C}_\omega$.
  \item Exact form of component metrics $d_{\mathrm{stance}}, d_{\mathrm{inst}}$ within the structure of $\delta$.
  \item Perturbation families used for robustness testing.
  \item Grounding weights $w_\tau$, saturation constants $K_\tau$, and mixing weight $\lambda$.
\end{itemize}

Any change to normative elements invalidates DRMM compliance.

\subsection*{A.10 \quad Forbidden Operations}

\paragraph{Human–readable.}
The safety spec is defined as much by what DRMM refuses to do as by what it does.

\paragraph{Canonical form.}
The following operations are forbidden in any system claiming DRMM compliance:
\begin{enumerate}[label=(F\arabic*)]
  \item Reporting a Level–0 vector, average, or trajectory as a ``concept'' or \texttt{concept\_vector}.
  \item Making a conceptual evolution claim for a branch that has not passed $\Pi_G$ and $\Pi_R$.
  \item Averaging or merging branches across any pair $(j,k)$ with $\delta_{jk} \ge \delta_{\mathrm{thr}}$.
  \item Dropping branches that fail gates without emitting the corresponding $\bot_\ast$ status.
  \item Omitting provenance fields defined in Section~\ref{sec:A8}, or misreporting threshold values.
  \item Comparing fields $\mathcal{F}_\omega$ and $\mathcal{F}_{\omega'}$ as if they lived in a shared state space without an explicitly defined inter–field mapping (to be specified outside this appendix).
\end{enumerate}

Any violation of these rules invalidates DRMM certification for the affected outputs.


```latex
\appendix

\section*{Appendix B \quad Operator-Theoretic Structure and Norm Bounds}

This appendix records explicit operator-theoretic properties and norm bounds that arise from the DRMM construction.
The purpose is twofold:
(i) to make clear the stability guarantees implied by the gate structure, and
(ii) to provide reusable lemmas and theorems that can be cited in later work without re-derivation.

Throughout, we work in finite-dimensional Hilbert spaces, so all norms are equivalent and all linear operators are bounded; we adopt the Euclidean norm on $\mathbb{R}^d$ and the associated operator norm on linear maps.

\section{Basic Normed Structure}

\subsection{Branch Space as a Product Hilbert Space}

Let
\[
\mathcal{H}_\mu = \mathbb{R}^d,
\quad
\mathcal{H}_\sigma = \mathbb{R}^{d_\sigma},
\quad
\mathcal{H}_\iota = \mathbb{R}^{d_\iota},
\quad
\mathcal{H}_\gamma = \mathbb{R},
\]
with the usual Euclidean inner products.
Define the branch state space as the direct product
\[
\mathcal{H}_\theta
  = \mathcal{H}_\mu \times \mathcal{H}_\sigma \times
    \mathcal{H}_\iota \times \mathcal{H}_\gamma.
\]

We equip $\mathcal{H}_\theta$ with the weighted norm
\[
\|\theta\|_{\Theta}
=
\sqrt{
  \alpha \|\mu\|_2^2
  + \beta \|\sigma\|_2^2
  + \gamma \|\iota\|_2^2
  + \eta \gamma^2
},
\]
for $\theta = (\mu,\sigma,\iota,\gamma)$ and fixed positive weights $\alpha,\beta,\gamma,\eta > 0$.
This defines a Hilbert-space structure equivalent to any other Euclidean norm on $\mathcal{H}_\theta$.

\subsection{Tension Metric as a Norm-Induced Distance}

Recall the tension metric
\[
\delta(\theta_i,\theta_j)
=
\alpha \,\lVert \mu_i-\mu_j \rVert_2
+
\beta \, d_{\mathrm{stance}}(\sigma_i,\sigma_j)
+
\gamma \, d_{\mathrm{inst}}(\iota_i,\iota_j)
+
\eta \, |\gamma_i-\gamma_j|.
\]

\paragraph{Assumption.}
We assume:
\begin{itemize}[nosep]
  \item $d_{\mathrm{stance}}$ and $d_{\mathrm{inst}}$ are metrics induced by norms on $\mathcal{H}_\sigma$ and $\mathcal{H}_\iota$, respectively:
  \[
  d_{\mathrm{stance}}(\sigma_i,\sigma_j) = \|\sigma_i-\sigma_j\|_{\Sigma},\quad
  d_{\mathrm{inst}}(\iota_i,\iota_j) = \|\iota_i-\iota_j\|_{I},
  \]
  where $\|\cdot\|_{\Sigma}$ and $\|\cdot\|_{I}$ are norms equivalent to the Euclidean norm.
\end{itemize}

\begin{lemma}[Norm Equivalence]
On $\mathcal{H}_\theta$, there exist constants $c_1,c_2 > 0$ such that for all $\theta_i,\theta_j$:
\[
c_1 \|\theta_i-\theta_j\|_{\Theta}
\;\le\;
\delta(\theta_i,\theta_j)
\;\le\;
c_2 \|\theta_i-\theta_j\|_{\Theta}.
\]
\end{lemma}

\begin{proof}[Sketch]
Because all component norms are equivalent in finite dimensions, there exist constants $k_\mu,k_\sigma,k_\iota,k_\gamma > 0$ such that
\[
\|\mu_i-\mu_j\|_2 \le k_\mu \|\mu_i-\mu_j\|_2,\quad
\|\sigma_i-\sigma_j\|_{\Sigma} \le k_\sigma \|\sigma_i-\sigma_j\|_2,
\]
and similarly for $\iota$ and $\gamma$.
The weighted sum defining $\delta$ is a continuous norm-equivalent functional of $(\mu_i-\mu_j,\sigma_i-\sigma_j,\iota_i-\iota_j,\gamma_i-\gamma_j)$.
Because all norms on a finite-dimensional vector space are equivalent, we obtain constants $c_1,c_2$ satisfying the inequality.
\end{proof}

\paragraph{Consequence.}
Any Lipschitz or contraction property stated in terms of $\delta$ can be restated in terms of the Hilbert norm $\|\cdot\|_{\Theta}$, and vice versa, up to constant factors.

\section{Gate Operators as Nonexpansive Maps}

In this section we analyze the gate operators $\Pi_G,\Pi_R,\Pi_D$ at the operator level, emphasizing their nonexpansive or contractive properties.

\subsection{Grounding Gate $\Pi_G$}

$\Pi_G$ maps branch states to either themselves or the rejection symbol $\bot_G$:
\[
\Pi_G(\theta_j) =
\begin{cases}
\theta_j & \gamma_j \ge \gamma_{\min},\\[2pt]
\bot_G   & \gamma_j < \gamma_{\min}.
\end{cases}
\]

We model $\bot_G$ as an absorbing non-Hilbert symbol external to $\mathcal{H}_\theta$.
Restricted to the admissible region
\[
\mathcal{H}_\theta^{(G)} = \{\theta \in \mathcal{H}_\theta : \gamma \ge \gamma_{\min}\},
\]
$\Pi_G$ acts as the identity.

\begin{proposition}[Nonexpansiveness of $\Pi_G$ on admissible region]
For $\theta_i,\theta_j \in \mathcal{H}_\theta^{(G)}$,
\[
\delta\bigl(\Pi_G(\theta_i),\Pi_G(\theta_j)\bigr)
=
\delta(\theta_i,\theta_j),
\]
and therefore $\Pi_G$ is an isometry and thus nonexpansive on $\mathcal{H}_\theta^{(G)}$.
\end{proposition}

\begin{proof}
On $\mathcal{H}_\theta^{(G)}$ we have $\Pi_G(\theta)=\theta$.
Thus the distance is unchanged.
\end{proof}

\paragraph{Remark.}
Extending $\delta$ to include $\bot_G$ requires choosing a finite or infinite penalty distance.
If we declare $\delta(\bot_G, \theta) = +\infty$ for $\theta\neq\bot_G$, then $\Pi_G$ is nonexpansive when both arguments remain admissible and expansive when one is rejected.
For norm bounds, we focus on the admissible region.

\subsection{Robustness Gate $\Pi_R$}

$\Pi_R$ depends on a family of perturbations $\{\theta_j^{(k)}\}_{k}$ and a branch-distance $d_{\mathrm{branch}}$:
\[
\Pi_R(\theta_j) =
\begin{cases}
\theta_j & \max_k d_{\mathrm{branch}}(\theta_j,\theta_j^{(k)}) < \epsilon_R,\\[2pt]
\bot_R   & \text{otherwise.}
\end{cases}
\]

Let $\Theta^{(R)}$ be the subset of branches that pass robustness for a fixed perturbation family; then $\Pi_R$ is the identity on $\Theta^{(R)}$ and acts as rejection otherwise.

\begin{proposition}[Nonexpansiveness on robust subset]
For $\theta_i,\theta_j \in \Theta^{(R)}$,
\[
\delta\bigl(\Pi_R(\theta_i),\Pi_R(\theta_j)\bigr)
=
\delta(\theta_i,\theta_j),
\]
so $\Pi_R$ is an isometry on the robust subset.
\end{proposition}

\begin{proof}
Identical to $\Pi_G$ on its admissible region.
\end{proof}

\paragraph{Note.}
The nontrivial property of $\Pi_R$ is not metric contraction but \emph{filtering}: it enforces that only branches stable across the declared perturbation family survive.

\subsection{Dialectical Gate $\Pi_D$ as Nonexpansive Selection}

$\Pi_D$ acts at the field level.
Given a set of branches $\{\theta_j\}_j$ and tensions $\delta_{jk}$, we:
\begin{itemize}[nosep]
  \item build a high-tension graph $H$ on indices $j$ with edges for pairs $(j,k)$ where $\delta_{jk} \ge \delta_{\mathrm{thr}}$,
  \item enforce that no two nodes connected by an edge in $H$ are merged,
  \item assign field-level status and maintain branch-level statuses.
\end{itemize}

$\Pi_D$ is not a linear operator on $\mathcal{H}_\theta$; it is a combinatorial operator on index sets.
However, it is \emph{nonexpansive} with respect to the maximum pairwise tension norm.

\begin{definition}[Field tension norm]
For a finite branch set $B=\{\theta_j\}_{j=1}^n$, define
\[
\|B\|_{\Delta} = \max_{i,j} \delta(\theta_i,\theta_j).
\]
\end{definition}

\begin{lemma}[Monotonicity of $\Pi_D$ under $\|\cdot\|_{\Delta}$]
Let $B$ be a branch set and $\Pi_D(B)$ the output branch set (after dialectical selection, without merging).
Then
\[
\|\Pi_D(B)\|_{\Delta} \le \|B\|_{\Delta}.
\]
\end{lemma}

\begin{proof}
$\Pi_D$ only removes branches or preserves them; it never introduces new branches or new pairwise distances.
Thus, the set of pairwise distances in $\Pi_D(B)$ is a subset of those in $B$, and the maximum cannot increase.
\end{proof}

\paragraph{Interpretation.}
In terms of the field tension norm, $\Pi_D$ is nonexpansive: it may reduce maximum tension by pruning branches but never amplifies it.

\section{Contraction and Stability of Iterative DRMM Steps}

Suppose we iterate DRMM across time or along an update sequence, e.g., with new data windows.
We wish to understand conditions under which the evolution of branch states is stable and admits fixed points.

\subsection{Single-Step Update as a Lipschitz Map}

Consider a deterministic update operator
\[
U : \mathcal{H}_\theta \to \mathcal{H}_\theta
\]
that maps branch states from time $t$ to $t+\Delta t$ before gate filtering:
\[
\theta_j(t+\Delta t) = U\bigl(\theta_j(t)\bigr).
\]

\paragraph{Assumption.}
Assume $U$ is Lipschitz with constant $L \ge 0$ with respect to $\delta$:
\[
\delta\bigl(U(\theta_i),U(\theta_j)\bigr) \le L \,\delta(\theta_i,\theta_j).
\]

\begin{lemma}[DRMM step as Lipschitz map]
Define the DRMM step map on a single branch as
\[
T = \Pi_G \circ \Pi_R \circ U.
\]
On the doubly admissible region
\[
\Theta^{(GR)} = \{\theta \in \mathcal{H}_\theta : \Pi_G(\theta) = \theta,\ \Pi_R(\theta) = \theta\},
\]
$T$ is Lipschitz with the same constant $L$:
\[
\delta\bigl(T(\theta_i),T(\theta_j)\bigr) \le L \,\delta(\theta_i,\theta_j)
\quad
\text{for all }
\theta_i,\theta_j \in \Theta^{(GR)}.
\]
\end{lemma}

\begin{proof}
On $\Theta^{(GR)}$, $\Pi_G$ and $\Pi_R$ act as identities.
Thus $T = U$ there, and the Lipschitz bound follows from the assumption on $U$.
\end{proof}

\subsection{Contraction and Fixed Points}

If $L < 1$, $T$ is a contraction mapping on the complete metric space $(\Theta^{(GR)},\delta)$ (assuming we restrict to a closed subset).

\begin{theorem}[Existence and uniqueness of stable branch state]
Suppose:
\begin{itemize}[nosep]
  \item $(\Theta^{(GR)},\delta)$ is complete and closed under $T$,
  \item $U$ is Lipschitz with constant $L<1$ on $\Theta^{(GR)}$.
\end{itemize}
Then:
\begin{enumerate}[nosep]
  \item $T$ has a unique fixed point $\theta^* \in \Theta^{(GR)}$,
  \item for any initial $\theta_0 \in \Theta^{(GR)}$, the iterates $\theta_{n+1}=T(\theta_n)$ converge to $\theta^*$.
\end{enumerate}
\end{theorem}

\begin{proof}
This is the Banach contraction mapping theorem applied to $T$ on the complete metric space $(\Theta^{(GR)},\delta)$.
\end{proof}

\paragraph{Interpretation.}
If the pre-gate update $U$ is contractive in tension, then the DRMM-filtered branch dynamics converge to a unique stable state.
This formalizes the intuitive notion of a ``settled'' conceptual regime for a branch.

\section{Norm Bounds for Grounding Coverage Updates}

Grounding coverage $\gamma_j$ is a scalar in $[0,1]$ depending on counts of grounding hooks.
Consider two branches $j$ and $k$; we can bound the contribution of grounding asymmetry to tension.

\begin{lemma}[Bound on grounding contribution]
For any branches $i,j$,
\[
0 \le \eta |\gamma_i - \gamma_j| \le \eta.
\]
\end{lemma}

\begin{proof}
Since $\gamma_i,\gamma_j \in [0,1]$, $|\gamma_i-\gamma_j| \le 1$.
Thus $\eta |\gamma_i-\gamma_j| \le \eta$.
\end{proof}

\paragraph{Consequence.}
The grounding asymmetry term adds a bounded offset to tension; all unbounded growth in $\delta$ must come from the geometric, stance, or institutional components.

\section{Field-Level Operator Norms}

Let $\mathcal{B}_n$ denote the set of branch sets of size at most $n$ for a fixed term $\omega$ and time window.
Equip $\mathcal{B}_n$ with the field tension norm
\[
\|B\|_{\Delta} = \max_{i,j} \delta(\theta_i,\theta_j).
\]

\subsection{DRMM Field Operator}

The DRMM field operator (single window) can be viewed as:
\[
\Phi : \mathcal{B}_n \to \mathcal{B}_n, \quad
\Phi(B) = \Pi_D \circ \Pi_R \circ \Pi_G \circ \mathcal{C}_\omega(B_0),
\]
where $B_0$ is an event-level representation and $\mathcal{C}_\omega$ produces branch candidates.

\paragraph{Assumptions.}
Assume:
\begin{itemize}[nosep]
  \item $\mathcal{C}_\omega$ is Lipschitz in the sense that small perturbations of events lead to small perturbations of branch states,
  \item the number of branches is uniformly bounded by $n$.
\end{itemize}

\begin{proposition}[Nonexpansiveness of field-level DRMM step]
Under the above assumptions and with fixed perturbation/regime,
there exists $L_\Phi \ge 0$ such that for all $B_1,B_2 \in \mathcal{B}_n$:
\[
\|\Phi(B_1) - \Phi(B_2)\|_{\Delta}
\le L_\Phi \|B_1 - B_2\|_{\Delta},
\]
with $L_\Phi$ determined by the Lipschitz constants of $\mathcal{C}_\omega$ and $U$, and the nonexpansiveness of $\Pi_G,\Pi_R,\Pi_D$.
In particular, if $\mathcal{C}_\omega$ and $U$ are nonexpansive, then so is $\Phi$.
\end{proposition}

\begin{proof}[Sketch]
$\Pi_G$ and $\Pi_R$ are isometries on their admissible regions; $\Pi_D$ is nonexpansive in $\|\cdot\|_{\Delta}$.
Thus, the Lipschitz constant of $\Phi$ is bounded by the product of the constants of $\mathcal{C}_\omega$ and $U$.
If both are nonexpansive ($L\le1$), so is $\Phi$.
\end{proof}

\section{Summary of Operator Norm Claims}

We summarize the key operator-norm-level claims established here:

\begin{enumerate}[label=(O\arabic*),nosep]
  \item \textbf{Norm equivalence.} The tension metric $\delta$ is equivalent to a Hilbert norm on $\mathcal{H}_\theta$, up to constant factors.
  \item \textbf{Gate isometries.} On their respective admissible regions, $\Pi_G$ and $\Pi_R$ are isometries, hence nonexpansive.
  \item \textbf{Dialectical nonexpansiveness.} $\Pi_D$ is nonexpansive with respect to the field tension norm $\|\cdot\|_{\Delta}$.
  \item \textbf{Lipschitz DRMM steps.} The single-branch DRMM update $T=\Pi_G\circ\Pi_R\circ U$ inherits the Lipschitz constant of $U$ on the doubly admissible region.
  \item \textbf{Contraction and fixed points.} If $U$ is a contraction on a closed, complete admissible region, then $T$ has a unique fixed point and iterates converge to it.
  \item \textbf{Grounding boundedness.} The grounding asymmetry term contributes at most $\eta$ to $\delta$ and is bounded.
  \item \textbf{Field-level nonexpansiveness.} The full DRMM field operator $\Phi$ is nonexpansive in $\|\cdot\|_{\Delta}$ when constructed from nonexpansive components.
\end{enumerate}

These claims together establish that DRMM is not only a semantic safety specification but also a stable operator on branch spaces and fields, with explicit norm bounds and convergence guarantees under reasonable assumptions.
They can be cited as prior art for any future work that:
\begin{itemize}[nosep]
  \item designs gate-based concept evolution operators on Hilbert-like branch spaces,
  \item uses tension metrics equivalent to Hilbert norms,
  \item and analyzes the stability of such operators via contraction mapping and nonexpansive operator theory.
\end{itemize}
```
\nocite{*}
\bibliographystyle{unsrt}  
\bibliography{references}  


\end{document}
